% ECONOMICS_PROBLEM: {"id": "C73-38", "title": "Classification and bifurcation of genuinely periodic MFG equilibria", "jel_code": "C73", "source_id": "OP-0217"}
% !TeX program = xelatex
\documentclass[11pt,a4paper]{article}
\usepackage[margin=23mm]{geometry}
\usepackage{fontspec}
\setmainfont{Latin Modern Roman}
\usepackage{amsmath,amssymb,mathtools}
\usepackage{xcolor,enumitem,booktabs,longtable,array,xurl,hyperref}
\definecolor{muted}{HTML}{53636C}
\hypersetup{colorlinks=true,urlcolor=blue,linkcolor=blue}
\setlength{\parindent}{0pt}
\setlength{\parskip}{5pt}
\newcommand{\R}{\mathbb R}
\newcommand{\E}{\mathbb E}
\newcommand{\N}{\mathbb N}
\newcommand{\ind}{\mathbf 1}
\newcommand{\Var}{\operatorname{Var}}
\newcommand{\Cov}{\operatorname{Cov}}
\newcommand{\supp}{\operatorname{supp}}
\DeclareMathOperator*{\argmin}{arg\,min}
\DeclareMathOperator*{\argmax}{arg\,max}
\newcommand{\entrymeta}[3]{{\sffamily\small\color{muted}\textbf{\texttt{#1}}\quad|\quad #2\par #3\par}}
\newcommand{\Problemref}[1]{\texttt{#1}}
\newcommand{\JELref}[2]{\texttt{#2} (JEL \texttt{#1})}
\begin{document}
\section*{Classification and bifurcation of genuinely periodic MFG equilibria}
\textbf{Problem ID:} \texttt{C73-38}\quad \textbf{JEL:} \texttt{C73}\par
% BEGIN ECONOMICS STATEMENT
\label{problem:OP-0217}
% GAME_THEORY_PROBLEM: {"local_id": "GT-087", "original_number": 87, "kind": "R", "entry_kind": "statement_only", "published": false, "review_required": true, "category": "game_theory"}
\label{game:GT-087}
{\sffamily\small\color{muted}
{\small\textbf{GT-087} | Mean-field games | R: a published agenda, with known existence examples\par}
\par}
\subsection*{Definitions and mathematical statement}
Fix $\nu>0$, an open parameter interval $I\subset\mathbb R$, and a smooth family of couplings $f_\eta:\mathbb T^d\times(0,\infty)\to\mathbb R$, $\eta\in I$. More generally a smooth uniformly convex Hamiltonian $H_\eta(x,p)$ may replace $|p|^2/2$; uniform convexity means $D_{pp}^2H_\eta\geq c\,\mathrm{Id}$ for a declared $c>0$.

Define $\mathcal P_f$ as the set of $\eta\in I$ for which some $P>0$, $\lambda\in\mathbb R$, and $C^{1,2}$ functions $(v,m)$ on $(\mathbb R/P\mathbb Z)\times\mathbb T^d$ satisfy
\[
 \begin{cases}
 \lambda-\partial_t v-\nu\Delta v+H_\eta(x,Dv)=f_\eta(x,m),\\
 \partial_t m-\nu\Delta m-
     \operatorname{div}\big(mD_pH_\eta(x,Dv)\big)=0,\\
 m>0,\quad\displaystyle\int_{\mathbb T^d}m(t,x)dx=1,
 \quad\displaystyle\int_0^P\int_{\mathbb T^d}v(t,x)dxdt=0,
 \end{cases}
\]
and $(Dv,m)$ is nonconstant in time. The actual value is $u=v-\lambda t$; a temporal drift in $u$ is compatible with periodic strategies and distributions. No fixed initial distribution or terminal cost is imposed in this periodic boundary-value problem.

Suppose a prescribed stationary branch $(\bar v_\eta,\bar m_\eta,\bar\lambda_\eta)$ exists. A periodic bifurcation point $(\eta_0,P_0)$ means there are nonconstant periodic solutions $(\eta_n,P_n,v_n,m_n,\lambda_n)$ with
\[
 \eta_n\to\eta_0,\quad P_n\to P_0,\quad
 (v_n(P_n\cdot),m_n(P_n\cdot),\lambda_n)
 \to(\bar v_{\eta_0},\bar m_{\eta_0},\bar\lambda_{\eta_0})
\]
in a declared $C^{1,2}$ topology on $(\mathbb R/\mathbb Z)\times\mathbb T^d$, with the spatial normalization of the stationary value chosen to match the periodic one.

\textbf{Research target.} Characterize $\mathcal P_f$ and its bifurcation points for declared families $(f_\eta,H_\eta)$, provide checkable sufficient/necessary hypotheses, and describe the extent of the resulting branches beyond a local neighborhood. This is a parameterized classification problem, not the assertion that every nonmonotone coupling has a periodic orbit.
\subsection*{Short English statement}
Which nonmonotone MFG systems admit genuine periodic population and strategy profiles, and how do the periodic branches emerge and continue?
\subsection*{Variants and scope boundaries}
Finite-horizon oscillations do not alone establish periodic solutions on the whole real line. First existence is already known: for the quadratic Hamiltonian and a spatially homogeneous local coupling, Cirant--Nurbekyan's Theorem~6.1 gives a specified bifurcation regime. General classification, global branch structure, and stability are further questions.
\subsection*{Source relationship and status}
The catalogue's source has a broad periodic-solutions discussion. The parameter set and bifurcation definition are editorial. Existing periodic examples must not be advertised as an unresolved first-existence question.
\subsection*{References}
\begingroup\footnotesize\raggedright
{\footnotesize\raggedright
\textbf{[S1]} Cardaliaguet and Delarue. \emph{Selected topics in mean field games}, ICM2022, vol.~5; DOI: 10.4171/ICM2022/17. \textit{Locator:} Section~3.3, ``Periodic solutions,'' p.~3680.\par
\url{https://ems.press/content/book-chapter-files/33265}\par\smallskip
\textbf{[S2]} Marco Cirant and Levon Nurbekyan (2018). \emph{The variational structure and time-periodic solutions for mean-field games systems}. Minimax Theory and its Applications 3, 227--260.\par
\textit{Locators:} Section~6, Theorem~6.1, pp.~24--25; Remark~6.4, p.~28.\par
\url{https://arxiv.org/pdf/1804.08943v1}\par\smallskip
\textbf{[S3]} Marco Cirant (2019). \emph{On the existence of oscillating solutions in non-monotone mean-field games}. Journal of Differential Equations 266, 8067--8093; DOI: 10.1016/j.jde.2018.12.025.\par
\textit{Locator:} arXiv version, introduction, p.~5: distinction between oscillations on a finite interval and truly periodic solutions, with a reference to the subsequent Cirant--Nurbekyan work.\par
\url{https://arxiv.org/pdf/1711.08047}\par}

% END AUDIT NOTES (not typeset):
% GT078: R; time-norm unspecified in source; centered values avoid secular errors.
% GT079: H for blanket monotone historical target, with 2025 resolution source.
% GT080--083: R; formal models/topologies/stability/noise class are editorial.
% GT084: P; finite-state potential, d>=3, fixed horizon and Markov selection.
% GT085: P only for unchanged running costs; corrected f^epsilon result is known.
% GT086: R nonpotential extension; weak potential uniqueness is known already.
% GT087: R classification; existence of actual periodic examples is known.
% Retrieval: EMS PDF equations/pages; Porretta--Rossi PDF; Cirant--Nurbekyan PDF;
% Cirant PDF; Cardaliaguet--Maillet--Yan 2025 PDF; Cecchin--Delarue abstract/metadata.
% No exhaustive post-source literature audit or independent proof audit is claimed.


% Entries 88--100.  Source versions were inspected on 8 October 2026.
% P: source-stated question; R: editorial research formulation;
% H: historical resolved/invalid formulation; U: verification limitation.
\par\endgroup

\subsection*{Common conventions: Source mathematical conventions}

\textbf{Finite games.} For a finite set $A$, $\Delta(A)$ is its probability
simplex. Players choose independent mixed strategies in Nash equilibrium;
correlation is allowed only when explicitly part of the solution concept.
In a game with payoff functions $u_i$ and strategy profile $x$, an additive
$\varepsilon$ Nash equilibrium satisfies
\[
 u_i(x)\geq u_i(x_i',x_{-i})-\varepsilon
 \quad\text{for every player $i$ and every admissible deviation $x_i'$}.
\]
For numerical approximation, payoffs are normalized as locally specified.
An $\varepsilon$ payoff residual is distinct from distance at most
$\varepsilon$ to the set of exact equilibria. Point convergence, convergence
of empirical play, and convergence in distance to an equilibrium set are
also distinct.

\textbf{Computational representations.} Unless stated otherwise, a complexity
question takes rational data with binary encoding; polynomial time is measured
in total bit length. A fixed-accuracy polynomial algorithm is not necessarily
polynomial in $1/\varepsilon$, and neither is necessarily polynomial in
$\log(1/\varepsilon)$. Succinct games and value, demand, or sampling oracles
must declare their interfaces. Exact real solutions require an explicit
representation; an unspecified list of arbitrary real numbers is not a
finite computational output. A claimed algorithmic barrier is meaningful
only for its specified model and reductions.

\textbf{Dynamic games.} Histories, information, observation, and allowable
randomization are part of the model. A stationary strategy depends only on
the current state; a positional strategy in a deterministic graph game
chooses one action at each controlled vertex. Nash equilibrium at an initial
state and equilibrium after every history are different requirements.
An expectation of a pathwise $\liminf$ payoff, a $\liminf$ of expected averages,
a limiting expected average, and a uniform finite-horizon equilibrium are
different criteria. None is silently substituted for another.

\textbf{Allocation and mechanisms.} A complete allocation assigns every item
exactly once unless otherwise stated. Zero-valued goods, negative values,
capacity constraints, feasibility, lotteries, and copies of goods affect
fairness definitions and are specified locally. Dominant-strategy incentive
compatibility, Bayesian incentive compatibility, universal truthfulness,
and truthfulness in expectation impose different inequalities. Whenever
an objective ratio has zero denominator, its convention is stated or those
instances are excluded.

\textbf{Infinite-dimensional models.} Expectations, integrals, stochastic
equations, and optimization objectives require their local regularity,
integrability, filtrations, and admissible domains. A backward--forward
mean-field system is a boundary-value problem; it is not an initial-value
semiflow without an additional construction. Long-time, large-population,
and vanishing-noise limits require an order of limits and a convergence
topology. If a minimum need not be attained, the objective is its infimum
or supremum and the approximation requirement is stated separately.
% END ECONOMICS STATEMENT
\end{document}
